\documentclass[11pt,oneside]{article}

\usepackage[T1]{fontenc}
\usepackage[utf8]{inputenc}
\usepackage{libertinus}
\usepackage[amsthm]{libertinust1math}
\usepackage{microtype}
\usepackage[margin=1.02in]{geometry}
\usepackage{mathtools,amssymb,amsthm,bm}
\usepackage{booktabs,longtable,array,tabularx}
\usepackage{enumitem}
\usepackage{xcolor}
\usepackage{graphicx}
\usepackage{url}
\usepackage[hidelinks]{hyperref}
\usepackage[nameinlink,noabbrev]{cleveref}
\usepackage{fancyhdr}
\usepackage{titlesec}
\usepackage{needspace}

\definecolor{oliveink}{RGB}{62,73,45}
\definecolor{softgray}{RGB}{90,90,90}
\hypersetup{
  colorlinks=true,
  linkcolor=oliveink,
  citecolor=oliveink,
  urlcolor=oliveink,
  pdftitle={Integral and Transform Calculus for the Fabius--Rvachev--Quantile System},
  pdfauthor={Research report prepared for Vladimir Reshetnikov},
  pdfsubject={Fabius, inverse Fabius, Rvachev up-function, integrals and transforms},
  pdfkeywords={Fabius function, Rvachev up-function, inverse Fabius function, integral transforms, order statistics, sinc products}
}

\pagestyle{fancy}
\fancyhf{}
\fancyhead[L]{\small Integral and Transform Calculus for the Fabius--Rvachev System}
\fancyhead[R]{\small\thepage}
\renewcommand{\headrulewidth}{0.35pt}
\setlength{\headheight}{14pt}

\titleformat{\section}{\Large\bfseries\color{oliveink}}{\thesection}{0.7em}{}
\titleformat{\subsection}{\large\bfseries}{\thesubsection}{0.65em}{}
\titleformat{\subsubsection}{\normalsize\bfseries}{\thesubsubsection}{0.6em}{}

\newtheorem{theorem}{Theorem}[section]
\newtheorem{proposition}[theorem]{Proposition}
\newtheorem{corollary}[theorem]{Corollary}
\newtheorem{lemma}[theorem]{Lemma}
\newtheorem{conjecture}[theorem]{Conjecture}
\newtheorem{problem}[theorem]{Problem}
\theoremstyle{definition}
\newtheorem{definition}[theorem]{Definition}
\newtheorem{algorithm}[theorem]{Algorithm}
\theoremstyle{remark}
\newtheorem{remark}[theorem]{Remark}
\newtheorem{example}[theorem]{Example}

\newcommand{\F}{F}
\newcommand{\G}{G}
\newcommand{\Up}{\operatorname{up}}
\newcommand{\E}{\mathbb E}
\newcommand{\Pp}{\mathbb P}
\newcommand{\R}{\mathbb R}
\newcommand{\C}{\mathbb C}
\newcommand{\N}{\mathbb N}
\newcommand{\dd}{\,\mathrm d}
\newcommand{\ii}{\mathrm i}
\newcommand{\sinc}{\operatorname{sinc}}
\newcommand{\Beta}{\operatorname{Beta}}
\newcommand{\PV}{\operatorname{PV}}
\newcommand{\Li}{\operatorname{Li}}
\newcommand{\fall}[1]{#1^{\underline{\vphantom{A}}}}
\newcommand{\tagnew}{\textsuperscript{\textcolor{oliveink}{\bfseries N}}}
\newcommand{\tagbase}{\textsuperscript{\textcolor{softgray}{\bfseries B}}}
\newcommand{\tagconj}{\textsuperscript{\textcolor{oliveink}{\bfseries C}}}
\newcommand{\Tn}{T_n}
\newcommand{\binomc}[2]{\binom{#1}{#2}}

\DeclareMathOperator{\Var}{Var}
\DeclareMathOperator{\Cov}{Cov}
\DeclareMathOperator{\supp}{supp}
\DeclareMathOperator{\sgn}{sgn}
\DeclareMathOperator{\RePart}{Re}
\DeclareMathOperator{\ImPart}{Im}

\setlist[itemize]{leftmargin=1.55em,itemsep=0.25em,topsep=0.35em}
\setlist[enumerate]{leftmargin=1.75em,itemsep=0.25em,topsep=0.35em}
\allowdisplaybreaks

\title{\bfseries Integral and Transform Calculus\\[0.25em]
for the Fabius--Rvachev--Quantile System\\[0.65em]
\large Order-statistic spacings, beta--quantile lattices, dyadic resolvents,\\
sinc energies, and beyond-all-orders localization}
\author{Research report prepared for Vladimir Reshetnikov}
\date{August 27, 2026}

\begin{document}
\maketitle

\begin{abstract}
This report develops a unified calculus for indefinite and definite integrals,
iterated and fractional antiderivatives, and Mellin, Laplace, Fourier,
Cauchy--Stieltjes, and beta transforms involving the bounded Fabius function
\(\F\), Rvachev's compactly supported function \(\Up\), and the inverse Fabius
function \(\G=\F^{-1}\).  The repository corpus was audited at the pinned
snapshot \texttt{a6555a64671b}; its recent
monomial/fractional-calculus, inverse-asymptotic, Fourier-decay, and
Bell--Bernoulli developments are treated as baseline rather than rediscovered.

The main new organizing principle is an adjacent-order-statistic spacing
identity.  It converts every weighted mixed integral
\(\int h(x)\F(x)^p(1-\F(x))^q\dd x\) into the expected increment of an
antiderivative of \(h\) across two consecutive order statistics.  This yields,
in one stroke, explicit formulas for polynomial, Jacobi, exponential,
trigonometric, logarithmic, Lambert-\(W\), Mellin, Fourier, and generalized
resolvent weights, as well as a closed probabilistic form for every integer or
fractional Volterra primitive.  A complementary beta--quantile calculus gives
incomplete antiderivatives in terms of truncated moments of \(\G\), a Pascal
lattice, reflection identities, and strict checkerboard total positivity.

Further results include: a monotone beta-normalized balanced sequence converging
to \(1/2\) faster than every inverse power; its logarithmic decay constant
\(-1/(8\log2)\), with a Lambert-\(W\) saddle; a Cauchy--Stieltjes hierarchy for
the \(\Up\)-law satisfying exact dyadic functional-differential equations; a
CDF/quantile resolvent duality; endpoint derivative growth of order
\(\exp((\log2)n^2/2)\); and an exact complementary pair of sinc-product energy
integrals for the previously opaque constant \(\int_0^1\F(x)^2\dd x\).
A commented Python program independently checks the principal identities and
is included with the source and PDF.
\end{abstract}

\tableofcontents
\newpage

\section{Scope, corpus audit, and status of the results}

\subsection{Audit boundary}

The audit covered the \texttt{*.tex} sources under
\begin{center}
\url{Analysis/FabiusFunction/docs}
\end{center}
including all subdirectories.  The working snapshot was pinned to commit
\texttt{a6555a64671b}, after the merge of the recent
frontier-report stack.  The pinned Git-tree enumeration used in the audit contained 87 TeX paths;
its complete path manifest is included with this report.  Byte-identical copies, source/PDF bundles, OCR transcriptions of external papers,
and successive draft variants were grouped by content lineage; theorem-level
comparison was performed against every content-distinct Fabius/Rvachev source
and against the archive ledger.

The closest predecessors are the reports on monomial antiderivatives,
complex-power Newton--Volterra expansions, fractional calculus, inverse-Fabius
asymptotics, derivative norm spectra, and the consolidated frontier reports.
Consequently, the following are treated here as \emph{baseline}:
\begin{itemize}
  \item the primitive ladder for \(x^\alpha\F(x)\), including the terminating
        integer-power case and convergent complex-power series;
  \item basic Volterra/Riemann--Liouville formulas for \(\F\) and \(\G\);
  \item endpoint asymptotics expressed through Lambert \(W_{-1}\), periodic
        corrections, Bell polynomials, and Gamma--zeta operators;
  \item the infinite sinc product for \(\widehat{\Up}\), its dyadic refinement,
        and the existing Fourier-decay analysis;
  \item exact dyadic values, rational moment recurrences, finite Thue--Morse
        splines, and the established \(q\)-binomial machinery.
\end{itemize}

A superscript \(\tagnew\) marks a result that appears new relative to the audited
repository corpus; \(\tagbase\) marks a recovered or reorganized baseline fact;
and \(\tagconj\) marks a conjecture.  These labels are not claims of global
priority in the mathematical literature.

\subsection{Principal new contributions}

The report's main contributions are:
\begin{enumerate}
  \item the weighted adjacent-spacing theorem and its partial, iterated, and
        fractional antiderivative forms;
  \item an incomplete beta--quantile antiderivative and a two-parameter
        beta--quantile lattice;
  \item a bivariate resolvent for all mixed powers, with Möbius symmetry and
        strict checkerboard total positivity;
  \item a beta-localization theorem, strict monotonicity of the balanced
        normalized sequence, superalgebraic convergence, and the exact leading
        logarithmic rate governed by a Lambert saddle;
  \item a Cauchy--Stieltjes hierarchy for \(\Up\), including boundary-value and
        endpoint nonanalyticity statements;
  \item a generalized resolvent duality interchanging \(\F\) and \(\G\);
  \item the sinc-energy identity and complementary Gini sum rule for
        \(\int_0^1\F^2\);
  \item a certified rational finite-spline approximation scheme inherited from
        the Thue--Morse truncations, with an explicit coupling error bound.
\end{enumerate}

\section{The three functions and their probability models}

\subsection{The bounded Fabius distribution}

Let \((U_j)_{j\ge1}\) be independent \(\operatorname{Unif}[0,1]\) variables and
set
\begin{equation}\label{eq:Xdef}
  X=\sum_{j\ge1}\frac{U_j}{2^j}.
\end{equation}
The bounded Fabius function is the distribution function
\begin{equation}
  \F(x)=\Pp(X\le x),\qquad 0\le x\le1,
\end{equation}
continued by \(0\) to the left and \(1\) to the right whenever a genuine CDF is
needed.  It obeys
\begin{equation}\label{eq:Fbasic}
 \F(1-x)=1-\F(x),\qquad
 \F'(x)=2\F(2x)\quad(0\le x\le\tfrac12),
\end{equation}
and is strictly increasing on \([0,1]\).  Its inverse is
\begin{equation}
  \G(u)=\F^{-1}(u),\qquad 0\le u\le1,
\end{equation}
with \(\G(1-u)=1-\G(u)\).

Write
\begin{equation}
  d_s=\E X^s=\int_0^1 x^s\F'(x)\dd x
\end{equation}
whenever the moment exists.  For integer \(n\ge1\), the exact recurrence is
\begin{equation}\label{eq:momentrec}
 (n+1)(2^n-1)d_n=\sum_{k=0}^{n-1}\binom{n+1}{k}d_k.
\end{equation}
In particular,
\[
 d_0=1,
 \quad d_1=\frac12,
 \quad d_2=\frac5{18},
 \quad d_3=\frac16,
 \quad d_4=\frac{143}{1350}.
\]

\subsection{Rvachev's compactly supported density}

Put
\begin{equation}\label{eq:Ydef}
 Y=2X-1=\sum_{j\ge1}\frac{V_j}{2^j},
 \qquad V_j=2U_j-1\sim\operatorname{Unif}[-1,1].
\end{equation}
Then \(Y\) has density
\begin{equation}\label{eq:updef}
 \Up(y)=\F(1-|y|)\,\mathbf1_{[-1,1]}(y),
 \qquad
 \F'(x)=2\Up(2x-1).
\end{equation}
Thus \(\Up\) is even and \(C^\infty\), and
\begin{equation}\label{eq:uprefinement}
 \Up'(x)=2\bigl(\Up(2x+1)-\Up(2x-1)\bigr).
\end{equation}
We use the cycles-per-unit Fourier convention
\begin{equation}
 \widehat f(\xi)=\int_{\R}f(x)e^{-2\pi\ii x\xi}\dd x,
 \qquad \sinc t=\frac{\sin t}{t}.
\end{equation}
Then
\begin{equation}\label{eq:phi}
 \Phi(\xi):=\widehat{\Up}(\xi)
 =\prod_{j\ge0}\sinc\!\left(\frac{\pi\xi}{2^j}\right).
\end{equation}
The moment-generating functions are
\begin{align}
 M_X(z)&=\E e^{zX}
 =\prod_{j\ge1}\frac{e^{z/2^j}-1}{z/2^j},\label{eq:MX}\\
 M_Y(z)&=\E e^{zY}
 =\prod_{j\ge1}\frac{\sinh(z/2^j)}{z/2^j}.\label{eq:MY}
\end{align}

\subsection{Bernoulli cumulants and Bell moments}

The transform \eqref{eq:MY} makes the Bell--Bernoulli structure transparent.
Since
\[
 \log\frac{\sinh z}{z}
 =\sum_{m\ge1}\frac{2^{2m-1}B_{2m}}{m(2m)!}z^{2m},
\]
one obtains
\begin{equation}\label{eq:cumulants}
 \log M_Y(z)
 =\sum_{m\ge1}
 \frac{2^{2m-1}B_{2m}}{m(2m)!(2^{2m}-1)}z^{2m}.
\end{equation}
Hence the nonzero cumulants are
\begin{equation}
 \kappa_{2m}(Y)=
 \frac{2^{2m-1}B_{2m}}{m(2^{2m}-1)},
 \qquad \kappa_{2m+1}(Y)=0,
\end{equation}
and the centered moments
\(c_m=\E Y^{2m}\) are complete Bell polynomials in these cumulants.  The first
values are
\[
 c_0=1,\quad c_1=\frac19,\quad c_2=\frac{19}{675},\quad
 c_3=\frac{583}{59535}.
\]
They also satisfy \(c_m=2d_{2m+1}/(2m+1)\).

\section{Endpoint integrability: powers of a CDF versus powers of a quantile}

The endpoint law established in the corpus has the logarithmic form
\begin{equation}\label{eq:endpointlog}
 \log\F(e^{-t})=-\frac{t^2}{2\log2}+O(t\log t),
 \qquad t\to+\infty,
\end{equation}
with much sharper Lambert-periodic expansions available in the baseline
reports.  This immediately creates a useful dichotomy.

\begin{proposition}[Integrability classification]\label{prop:integrability}
Let \(p,q,\alpha,\beta\in\C\).
\begin{enumerate}[label=(\roman*)]
 \item The mixed power
 \(\F(x)^p(1-\F(x))^q\) is absolutely integrable on \((0,1)\) precisely when
 \(\Re p\ge0\) and \(\Re q\ge0\).
 \item If \(\Re p>0\), then \(x^\alpha\F(x)^p\) is integrable at \(0\) for
 every \(\alpha\).  If \(\Re p=0\), one needs \(\Re\alpha>-1\); if
 \(\Re p<0\), no algebraic weight repairs the divergence.
 \item Every fixed negative moment of \(X\) is finite:
 \[
   \E X^{-s}=\int_0^1\G(u)^{-s}\dd u<\infty\qquad(s>0),
 \]
 whereas \(\int_0^1\F(x)^{-s}\dd x=\infty\).
 \item For \(\Re\alpha>-1\), the integral
 \(\int_{-1}^1|x|^\alpha\Up(x)^q\dd x\) converges when \(\Re q\ge0\) and
 diverges when \(\Re q<0\).
\end{enumerate}
\end{proposition}

\begin{proof}
Put \(x=e^{-t}\).  By \eqref{eq:endpointlog}, a negative power
\(\F(x)^{-s}\) contributes
\(\exp(s t^2/(2\log2)-t+O(t\log t))\dd t\), which diverges.  A positive power
has the opposite quadratic sign and dominates every algebraic singularity.
For the quantile, inversion of \eqref{eq:endpointlog} gives
\(
 \G(e^{-L})=\exp(-\sqrt{2\log2\,L}+O(\log L))
\).
Thus \(\G(u)^{-s}\dd u\), under \(u=e^{-L}\), has exponent
\(-L+O(\sqrt L)\), and is integrable.  The remaining endpoint follows from
symmetry, while \(\Up(1-x)=\F(x)\) on \([0,1]\).
\end{proof}

This contrast explains why negative powers of \(\F\) or \(\Up\) generally do
not possess ordinary definite integrals, while inverse-Fabius negative moments
are perfectly well behaved.

\section{The adjacent-order-statistic spacing principle}

\subsection{The universal weighted identity}

Let \(Z_1,\ldots,Z_N\) be independent copies of \(X\), and write
\(Z_{1:N}\le\cdots\le Z_{N:N}\) for their order statistics.

\begin{theorem}[Weighted adjacent-spacing theorem\tagnew]\label{thm:spacing}
Let \(p,q\in\N_{>0}\), \(N=p+q\), and let \(h\in L^1(0,1)\) have an absolutely
continuous primitive \(H\), so \(H'=h\) almost everywhere.  Then
\begin{equation}\label{eq:spacing}
 \boxed{
 \int_0^1h(x)\F(x)^p(1-\F(x))^q\dd x
 =\frac{\E\bigl[H(Z_{p+1:N})-H(Z_{p:N})\bigr]}
        {\binom{N}{p}}.}
\end{equation}
More generally, for \(0\le a<b\le1\),
\begin{equation}\label{eq:spacinginterval}
 \int_a^b h(x)\F(x)^p(1-\F(x))^q\dd x
 =\E\!\left[
 \mathbf1_{\{L<U\}}
 \int_a^b h(x)\mathbf1_{[L,U)}(x)\dd x\right],
\end{equation}
where \(L\) is the maximum of \(p\) independent copies of \(X\), and \(U\)
is the minimum of another \(q\) copies.
\end{theorem}

\begin{proof}
For every fixed \(x\), independence gives
\[
 \Pp(L\le x<U)=\F(x)^p(1-\F(x))^q.
\]
Fubini's theorem yields \eqref{eq:spacinginterval}.  On \(\{L<U\}\), all
\(p\) observations in the first group precede all \(q\) observations in the
second.  Conditional on the unlabeled ordered sample, the labels are uniformly
distributed among the \(\binom Np\) possibilities.  Exactly one labeling has
the first \(p\) observations in the lower group.  On that labeling,
\(L=Z_{p:N}\) and \(U=Z_{p+1:N}\).  Integrating \(h\) across the random gap
gives \eqref{eq:spacing}.
\end{proof}

The theorem is distribution-free: the Fabius law enters through the order
statistics, not through the proof.  Its value here is that the order statistics
inherit the exact dyadic, inverse-quantile, and endpoint structures of \(\F\).

\subsection{Moments, Jacobi weights, exponentials, and Lambert weights}

Taking \(H(x)=x^{m+1}/(m+1)\) gives
\begin{equation}\label{eq:weightedmomspacing}
 \int_0^1x^m\F(x)^p(1-\F(x))^q\dd x
 =\frac{\E\bigl[Z_{p+1:N}^{m+1}-Z_{p:N}^{m+1}\bigr]}
 {(m+1)\binom Np}.
\end{equation}
For Jacobi weights, with the incomplete beta function
\(B_x(a,b)=\int_0^xt^{a-1}(1-t)^{b-1}\dd t\),
\begin{equation}\label{eq:jacobispacing}
 \int_0^1x^\alpha(1-x)^\beta\F(x)^p(1-\F(x))^q\dd x
 =\frac{\E\bigl[B_{Z_{p+1:N}}(\alpha+1,\beta+1)
                    -B_{Z_{p:N}}(\alpha+1,\beta+1)\bigr]}
 {\binom Np}.
\end{equation}
Differentiating in \(\alpha\) or \(\beta\) supplies all logarithmic Jacobi
weights.

For \(z\ne0\),
\begin{align}
 \int_0^1e^{zx}\F(x)^p(1-\F(x))^q\dd x
 &=\frac{\E(e^{zZ_{p+1:N}}-e^{zZ_{p:N}})}
 {z\binom Np},\label{eq:expspacing}\\
 \int_0^1e^{-2\pi\ii\xi x}\F(x)^p(1-\F(x))^q\dd x
 &=\frac{\E(e^{-2\pi\ii\xi Z_{p:N}}-e^{-2\pi\ii\xi Z_{p+1:N}})}
 {2\pi\ii\xi\binom Np}.\label{eq:fourierspacing}
\end{align}
The limits at \(z=0\) or \(\xi=0\) recover the mean spacing.

A Lambert-\(W\) weight also has an elementary primitive:
\begin{equation}
 \frac{\dd}{\dd x}
 \left(W(cx)+\frac12W(cx)^2\right)=\frac{W(cx)}{x}.
\end{equation}
Therefore, on any branch for which the path avoids the branch cut,
\begin{equation}\label{eq:lambertweight}
 \int_0^1\frac{W(cx)}x\F(x)^p(1-\F(x))^q\dd x
 =\frac{\E[\mathcal W_c(Z_{p+1:N})-\mathcal W_c(Z_{p:N})]}
 {\binom Np},
\end{equation}
where \(\mathcal W_c(x)=W(cx)+\tfrac12W(cx)^2\).  This is a direct integral
appearance of Lambert \(W\), distinct from its asymptotic role at the endpoint.

\subsection{Partial, iterated, and fractional antiderivatives}

For \(x\in[0,1]\), the partial primitive has the exact random-interval form
\begin{equation}\label{eq:partialspacing}
 \int_0^x h(t)\F(t)^p(1-\F(t))^q\dd t
 =\E\!\left[
 \mathbf1_{\{L<x,\,L<U\}}
 \bigl(H(\min\{x,U\})-H(L)\bigr)
 \right].
\end{equation}
This is already an indefinite integral: changing the lower endpoint adds only a
constant.

The formula becomes especially sharp for the Riemann--Liouville operator
\begin{equation}
 (I^\nu f)(x)=\frac1{\Gamma(\nu)}\int_0^x(x-t)^{\nu-1}f(t)\dd t,
 \qquad \Re\nu>0.
\end{equation}

\begin{theorem}[Fractional adjacent-spacing primitive\tagnew]\label{thm:fracspacing}
For \(p,q\in\N_{>0}\), \(N=p+q\), \(x\in[0,1]\), and \(\Re\nu>0\),
\begin{equation}\label{eq:fracspacing}
 \boxed{
 I^\nu\!\left[\F^p(1-\F)^q\right](x)
 =\frac{\E\left[(x-Z_{p:N})_+^\nu-(x-Z_{p+1:N})_+^\nu\right]}
 {\Gamma(\nu+1)\binom Np}.}
\end{equation}
For integer \(n\ge1\), the right side is the normalized difference of two
truncated \(n\)-th moments and is the unique \(n\)-fold primitive whose first
\(n\) initial data vanish at \(0\).
\end{theorem}

\begin{proof}
Insert
\(\F(t)^p(1-\F(t))^q=\Pp(L\le t<U)\) into the Volterra integral and apply
Fubini.  For fixed \(L<U\),
\[
 \frac1{\Gamma(\nu)}\int_0^x(x-t)^{\nu-1}\mathbf1_{[L,U)}(t)\dd t
 =\frac{(x-L)_+^\nu-(x-U)_+^\nu}{\Gamma(\nu+1)}.
\]
The same random-label argument as in \cref{thm:spacing} supplies the factor
\(\binom Np^{-1}\).
\end{proof}

Endpoint cases are also useful.  If \(M_p=Z_{p:p}\) and \(m_q=Z_{1:q}\), then
\begin{align}
 I^\nu[\F^p](x)
 &=\frac{\E(x-M_p)_+^\nu}{\Gamma(\nu+1)},\label{eq:Fpowerfrac}\\
 I^\nu[(1-\F)^q](x)
 &=\frac{x^\nu-\E(x-m_q)_+^\nu}{\Gamma(\nu+1)}.\label{eq:survfrac}
\end{align}
For \(p=1\), \eqref{eq:Fpowerfrac} collapses to the baseline Fabius ladder
\begin{equation}\label{eq:Fvolterraladder}
 \E(x-X)_+^n=n!\,2^{\binom n2}\F(x/2^n).
\end{equation}
Thus \cref{thm:fracspacing} is a genuine mixed-power extension of that ladder.

\section{Incomplete beta--quantile antiderivatives}

\subsection{A closed incomplete formula}

For \(r\in\C\), \(a,b>0\), and \(0\le y\le1\), define
\begin{equation}\label{eq:Qdef}
 Q_r(a,b;y)=\int_0^y\G(u)^r u^{a-1}(1-u)^{b-1}\dd u,
 \qquad Q_r(a,b)=Q_r(a,b;1).
\end{equation}

\begin{theorem}[Incomplete beta--quantile primitive\tagnew]\label{thm:incompleteQ}
Let \(m>-1\), \(p,q>0\), \(a=\F(x)\).  Whenever the displayed integrals
converge,
\begin{align}
 \int_0^x t^m\F(t)^p(1-\F(t))^q\dd t
 &=\frac{x^{m+1}a^p(1-a)^q}{m+1}\notag\\
 &\quad-\frac{p}{m+1}Q_{m+1}(p,q+1;a)
 +\frac{q}{m+1}Q_{m+1}(p+1,q;a).
 \label{eq:incompleteQ}
\end{align}
Consequently,
\begin{equation}\label{eq:defQ}
 \int_0^1 t^m\F(t)^p(1-\F(t))^q\dd t
 =\frac{qQ_{m+1}(p+1,q)-pQ_{m+1}(p,q+1)}{m+1}.
\end{equation}
\end{theorem}

\begin{proof}
Integrate by parts with \(t^m\dd t=\dd(t^{m+1}/(m+1))\).  The derivative of
\(\F^p(1-\F)^q\) is
\[
 \left[p\F^{p-1}(1-\F)^q-q\F^p(1-\F)^{q-1}\right]\F'.
\]
Substitute \(u=\F(t)\) in the two resulting terms.  The endpoint term vanishes
at \(x=1\) when \(q>0\).
\end{proof}

Formula \eqref{eq:incompleteQ} is a deterministic counterpart of the random-gap
formula \eqref{eq:partialspacing}.  It is often better for symbolic work because
all dependence on the inverse Fabius function is isolated in \(Q_r\).

\subsection{The beta--quantile lattice}

The complete transforms obey exact algebraic identities.

\begin{proposition}[Pascal, reflection, and order-statistic laws\tagnew]
For admissible parameters,
\begin{align}
 Q_r(a,b)&=Q_r(a+1,b)+Q_r(a,b+1),\label{eq:Qpascal}\\
 Q_1(a,b)+Q_1(b,a)&=B(a,b),\label{eq:Qreflection}\\
 Q_r(b,a)&=\sum_{k=0}^r(-1)^k\binom rk Q_k(a,b)
 \qquad(r\in\N).\label{eq:Qtriangular}
\end{align}
If \(a,b\in\N_{>0}\), then
\begin{equation}\label{eq:Qorder}
 Q_r(a,b)=B(a,b)\E Z_{a:a+b-1}^r.
\end{equation}
\end{proposition}

\begin{proof}
Equation \eqref{eq:Qpascal} is the identity
\(u^{a-1}(1-u)^{b-1}=u^a(1-u)^{b-1}+u^{a-1}(1-u)^b\).
For reflection, substitute \(u\mapsto1-u\) and use
\(\G(1-u)=1-\G(u)\).  The same substitution followed by the binomial theorem
gives \eqref{eq:Qtriangular}.  Finally,
\(\F(Z_{a:a+b-1})\sim\Beta(a,b)\), proving \eqref{eq:Qorder}.
\end{proof}

For \(q_s=Q_1(s,1)\), integration by parts gives
\begin{equation}\label{eq:qA}
 q_s=\int_0^1u^{s-1}\G(u)\dd u=\frac{1-A_s}{s},
 \qquad A_s:=\int_0^1\F(x)^s\dd x.
\end{equation}
When \(b\in\N\),
\begin{equation}\label{eq:Qfinitediff}
 Q_1(a,b)=\sum_{j=0}^{b-1}(-1)^j\binom{b-1}{j}q_{a+j}.
\end{equation}
Thus the entire integer beta--quantile lattice is a finite binomial transform of
the one-dimensional sequence \((A_s)\).

\subsection{Unaligned products of \(\F\) and \(\G\)}

Compositions aligned with the inverse are immediate:
\[
 \G(\F(x))=x,
 \qquad \F(\G(u))=u.
\]
Products at the same external argument, such as \(\F(x)^a\G(x)^b\), do not
collapse.  They nevertheless admit an exact two-dimensional layer-cake kernel.
For \(a,b>0\),
\begin{equation}\label{eq:FGkernel}
 \int_0^1\F(x)^a\G(x)^b\dd x
 =ab\int_0^1\!\int_0^1
 u^{a-1}v^{b-1}
 \bigl[1-\max\{\G(u),\F(v)\}\bigr]\dd u\dd v.
\end{equation}
Indeed,
\(\F(x)^a=a\int_0^1u^{a-1}\mathbf1_{\{\G(u)\le x\}}\dd u\), while
\(\G(x)^b=b\int_0^1v^{b-1}\mathbf1_{\{\F(v)\le x\}}\dd v\).  Formula
\eqref{eq:FGkernel} is a promising starting point for genuinely nonlinear
\(\F\)--\(\G\) coupling constants.

\section{The mixed-power lattice}

Define
\begin{equation}\label{eq:Ipq}
 I_{p,q}=\int_0^1\F(x)^p(1-\F(x))^q\dd x,
 \qquad p,q\ge0.
\end{equation}
For integer arguments, the elementary partition
\(1=\F+(1-\F)\) gives
\begin{equation}\label{eq:PascalI}
 I_{p,q}=I_{p+1,q}+I_{p,q+1},
 \qquad I_{p,q}=I_{q,p}.
\end{equation}
If \(\Delta A_p=A_{p+1}-A_p\), then
\begin{equation}\label{eq:finiteDI}
 I_{p,q}=(-1)^q\Delta^q A_p
 =\sum_{j=0}^q(-1)^j\binom qj A_{p+j}.
\end{equation}
For noninteger \(q\), an Abel-regularized Newton series remains valid:
\begin{equation}\label{eq:abelseries}
 I_{p,q}=\lim_{r\uparrow1}
 \sum_{j\ge0}(-1)^j\binom qj r^j A_{p+j},
 \qquad \Re p,\Re q\ge0.
\end{equation}

\subsection{The bivariate resolvent}

Let
\begin{equation}\label{eq:Rdef}
 R(z)=\int_0^1\frac{\dd x}{1-z\F(x)}
 =\sum_{p\ge0}A_pz^p,
 \qquad |z|<1.
\end{equation}

\begin{theorem}[Bivariate mixed-power generating function\tagnew]
For \(|z|,|w|<1\), with removable points understood by analytic continuation,
\begin{equation}\label{eq:bivgen}
 \boxed{
 \sum_{p,q\ge0}I_{p,q}z^pw^q
 =\frac{zR(z)+wR(w)}{z+w-zw}.}
\end{equation}
Moreover,
\begin{equation}\label{eq:Rmobius}
 R(z)=\frac1{1-z}R\!\left(-\frac z{1-z}\right).
\end{equation}
\end{theorem}

\begin{proof}
Sum the geometric series under the integral:
\[
 \sum_{p,q\ge0}I_{p,q}z^pw^q
 =\int_0^1\frac{\dd x}
 {(1-z\F(x))(1-w(1-\F(x)))}.
\]
Partial fractions give coefficients
\(z/(z+w-zw)\) and \(w/(z+w-zw)\).  Reflection
\(x\mapsto1-x\) changes the second integral into \(R(w)\).  The same reflection
applied directly to \(R\) gives \eqref{eq:Rmobius}.
\end{proof}

Parameter differentiation supplies mixed logarithmic integrals:
\begin{equation}\label{eq:paramdiff}
 \frac{\partial^{r+s}}{\partial p^r\partial q^s}I_{p,q}
 =\int_0^1\F(x)^p(1-\F(x))^q
 (\log\F(x))^r(\log(1-\F(x)))^s\dd x.
\end{equation}

\subsection{Strict checkerboard total positivity}

The lattice has much stronger structure than Pascal recursion.

\begin{theorem}[Strict sign-regularity\tagnew]\label{thm:signregular}
Let \(p_1<\cdots<p_m\) and \(q_1<\cdots<q_m\) be nonnegative integers.  Then
\begin{equation}\label{eq:signregular}
 (-1)^{\binom m2}
 \det\bigl[I_{p_i,q_j}\bigr]_{i,j=1}^m>0.
\end{equation}
Equivalently, reversing the order of the \(q\)-columns produces a strictly
totally positive matrix.
\end{theorem}

\begin{proof}
Under \(u=\F(x)\),
\[
 I_{p,q}=\int_0^1u^p(1-u)^q\G'(u)\dd u.
\]
The continuous Cauchy--Binet (Andreief) identity gives
\begin{align*}
 \det[I_{p_i,q_j}]
 &=\frac1{m!}\int_{(0,1)^m}
 \det[u_k^{p_i}]_{i,k}
 \det[(1-u_k)^{q_j}]_{j,k}
 \prod_{k=1}^m\G'(u_k)\dd u_k.
\end{align*}
On \(0<u_1<\cdots<u_m<1\), the first generalized Vandermonde determinant is
positive.  The numbers \(1-u_k\) occur in decreasing order, so the second has
sign \((-1)^{\binom m2}\).  Both are nonzero away from diagonals, and
\(\G'>0\), proving strictness.
\end{proof}

For \(m=2\), this yields the sharp reverse-TP inequality
\begin{equation}\label{eq:TP2}
 I_{p,q}I_{p+1,q+1}
 <I_{p,q+1}I_{p+1,q}.
\end{equation}
For fixed \(q\), the sequence \(p\mapsto I_{p,q}\) is a strict Hausdorff moment
sequence and hence has positive Hankel minors; \cref{thm:signregular} adds the
cross-parameter sign law.

\section{Beta localization and asymptotics}

\subsection{Interior localization}

Substitution \(u=\F(x)\) gives, for a weight \(h\),
\begin{equation}\label{eq:betaloc}
 \int_0^1h(x)\F(x)^p(1-\F(x))^q\dd x
 =B(p+1,q+1)\E K(B_{p,q}),
\end{equation}
where
\begin{equation}
 K(u)=h(\G(u))\G'(u),
 \qquad B_{p,q}\sim\Beta(p+1,q+1).
\end{equation}
Let \(N=p+q+2\) and \(\mu=(p+1)/N\).  If \(p/N\) remains in a compact
subinterval of \((0,1)\) and \(K\in C^4\) near \(\mu\), then
\begin{equation}\label{eq:betaexpansion}
 \E K(B_{p,q})
 =K(\mu)+\frac{\mu(1-\mu)}{2(N+1)}K''(\mu)+O(N^{-2}).
\end{equation}
In particular, when \(p+q=n\) and \(p/n\to\lambda\in(0,1)\),
\begin{equation}\label{eq:Ipqasymptotic}
 I_{p,q}
 \sim \sqrt{\frac{2\pi\lambda(1-\lambda)}{n}}
 \lambda^p(1-\lambda)^q\G'(\lambda).
\end{equation}
This is the ordinary, algebraic interior regime.

\subsection{The balanced sequence and central flatness}

The balanced case is exceptional.  Define
\begin{equation}\label{eq:Km}
 K_m=(2m+1)\binom{2m}{m}I_{m,m}
 =\E\G'(B_m),
 \qquad B_m\sim\Beta(m+1,m+1).
\end{equation}
The center identity
\begin{equation}\label{eq:centerF}
 \F\!\left(\frac12+h\right)
 =\frac12+2h-\sgn(h)\F(|h|)
\end{equation}
implies
\begin{equation}\label{eq:centerderivative}
 \F'\!\left(\frac12+h\right)=2[1-\F(2|h|)].
\end{equation}
Because every derivative of \(\F\) vanishes at \(0\), inverse differentiation
gives
\begin{equation}\label{eq:Gflat}
 \G'(1/2)=\frac12,
 \qquad \G^{(r)}(1/2)=0\quad(r\ge2).
\end{equation}
Thus \(\G'\) is flat, but not locally constant, at its minimum.

\begin{theorem}[Monotone superalgebraic localization\tagnew]\label{thm:Km}
The sequence \((K_m)_{m\ge0}\) is strictly decreasing and
\begin{equation}
 K_m\downarrow\frac12.
\end{equation}
Moreover, for every \(A>0\),
\begin{equation}\label{eq:superalg}
 K_m-\frac12=O_A(m^{-A}).
\end{equation}
\end{theorem}

\begin{proof}
The function \(\G'\) is symmetric about \(1/2\).  On \((1/2,1)\),
\(\F'\) is strictly decreasing by \eqref{eq:centerderivative}, so
\(\G'(u)=1/\F'(\G(u))\) is strictly increasing.  Hence \(\G'\) increases with
\(|u-1/2|\).  The ratio of successive beta densities is a positive constant
times \(u(1-u)\), which strictly decreases with \(|u-1/2|\).  A one-crossing
argument therefore gives \(\E\G'(B_{m+1})<\E\G'(B_m)\).  Concentration gives the
limit \(\G'(1/2)=1/2\).

For any \(A\), flatness \eqref{eq:Gflat} implies
\(|\G'(u)-1/2|\le C_A|u-1/2|^{2A}\) in a central neighborhood.  The corresponding
beta moment is \(O(m^{-A})\), while the complement has exponentially small beta
probability.  This proves \eqref{eq:superalg}.
\end{proof}

\subsection{A logarithmic law and Lambert saddle}

Superalgebraic decay conceals a precise logarithmic scale.

\begin{theorem}[Balanced logarithmic asymptotic\tagnew]\label{thm:Klog}
Assume the endpoint logarithm \eqref{eq:endpointlog}.  Then
\begin{equation}\label{eq:Klog}
 \boxed{
 \lim_{m\to\infty}
 \frac{\log(K_m-\tfrac12)}{(\log m)^2}
 =-\frac1{8\log2}.}
\end{equation}
The saddle in the central displacement \(v=|u-1/2|=e^{-t}\) is located at
\begin{equation}\label{eq:LambertSaddle}
 t_m=\frac12W(16m\log2)+o(1).
\end{equation}
\end{theorem}

\begin{proof}[Proof at logarithmic precision]
For \(u=1/2+v\), write \(\G(u)=1/2+h\).  Equations
\eqref{eq:centerF}--\eqref{eq:centerderivative} imply \(v=2h-\F(h)\sim2h\) and
\[
 \G'(1/2+v)-\frac12
 =\frac{\F(2h)}{2(1-\F(2h))}.
\]
Consequently,
\begin{equation}\label{eq:gprimeendpoint}
 \log\bigl(\G'(1/2+v)-\tfrac12\bigr)
 =-\frac{(\log(1/v))^2}{2\log2}
 +O(\log(1/v)\log\log(1/v)).
\end{equation}
The beta density contributes
\((1-4v^2)^m=\exp(-4mv^2+O(mv^4))\).  At logarithmic precision the relevant
exponent is therefore
\[
 \Psi_m(t)=4me^{-2t}+\frac{t^2}{2\log2}.
\]
Its stationary equation is
\(8me^{-2t}=t/\log2\), or
\((2t)e^{2t}=16m\log2\), giving \eqref{eq:LambertSaddle}.  Since
\(t_m\sim\tfrac12\log m\),
\[
 \Psi_m(t_m)=\frac{t_m^2+t_m}{2\log2}
 =\frac{(\log m)^2}{8\log2}+o((\log m)^2).
\]
Standard upper and lower Laplace bounds on windows around the saddle prove
\eqref{eq:Klog}.
\end{proof}

\begin{conjecture}[Full Lambert-periodic expansion\tagconj]\label{conj:Kfull}
Let \(w_m=W(16m\log2)\).  The quantity \(K_m-1/2\) has a full expansion of the
form
\begin{equation}\label{eq:Kfullconj}
 K_m-\frac12
 \sim \mathcal A(\vartheta_m)
 m^{\delta}w_m^{\gamma}\exp\!\left[-(w_m^2+2w_m)/(8\log2)\right]
 \left(1+\sum_{j\ge1}\frac{P_j(w_m)}{w_m^j}\right),
\end{equation}
where \(\mathcal A\) is a nonconstant bounded periodic factor inherited from the
endpoint phase, \(\vartheta_m\) is an affine function of
\(\log_2 w_m\), and the algebraic prefactor and polynomials \(P_j\) are determined by the
all-orders endpoint expansion of \(\F\).
\end{conjecture}

The exact periodic phase should be extracted from the corpus's Lambert-phase
normal form rather than guessed from moderate values of \(m\); the convergence
to the scale in \eqref{eq:Klog} is extremely slow.

\section{Powers of Rvachev's \(\Up\) function}

Because \(\Up(x)=1-\F(x)\) on \([0,1]\) and \(\Up\) is even, powers of \(\Up\)
reduce to inverse-beta moments.

\begin{theorem}[Power and weighted-power formulas\tagnew]\label{thm:uppowers}
For \(\Re\alpha>-1\) and \(\Re q>0\),
\begin{equation}\label{eq:uppower}
 \boxed{
 \int_{-1}^1|x|^\alpha\Up(x)^q\dd x
 =\frac{2q}{\alpha+1}
 \int_0^1\G(u)^{\alpha+1}(1-u)^{q-1}\dd u.}
\end{equation}
If \(q\in\N_{>0}\), then
\begin{equation}\label{eq:upmin}
 \int_{-1}^1|x|^\alpha\Up(x)^q\dd x
 =\frac2{\alpha+1}\E Z_{1:q}^{\alpha+1}.
\end{equation}
In particular,
\begin{equation}\label{eq:upLp}
 \int_{-1}^1\Up(x)^q\dd x=2A_q=2\E Z_{1:q}.
\end{equation}
\end{theorem}

\begin{proof}
On \([0,1]\), integrate \(x^\alpha(1-\F(x))^q\) by parts.  The boundary terms
vanish and substitution \(u=\F(x)\) gives
\[
 \int_0^1x^\alpha(1-\F(x))^q\dd x
 =\frac{q}{\alpha+1}\int_0^1\G(u)^{\alpha+1}(1-u)^{q-1}\dd u.
\]
Double by evenness.  For integer \(q\), the beta variable
\(\Beta(1,q)\) is \(\F(Z_{1:q})\).
\end{proof}

The incomplete one-sided primitive is
\begin{equation}\label{eq:upincomplete}
 \int_0^x t^\alpha\Up(t)^q\dd t
 =\frac{x^{\alpha+1}(1-\F(x))^q}{\alpha+1}
 +\frac{q}{\alpha+1}Q_{\alpha+1}(1,q;\F(x)).
\end{equation}
For \(q=1\),
\begin{equation}
 \int_{-1}^1|x|^\alpha\Up(x)\dd x
 =\frac{2d_{\alpha+1}}{\alpha+1}.
\end{equation}

\subsection{Coupling to the derivative of \(\F\)}

The derivative removes all distribution-specific complexity:
\begin{align}
 \int_0^1\F(x)^p(1-\F(x))^q\F'(x)\dd x
 &=B(p+1,q+1),\label{eq:Fprimebeta}\\
 \int_0^1\F(x)^p(1-\F(x))^q\Up(2x-1)\dd x
 &=\frac12B(p+1,q+1),\label{eq:upcoupling}\\
 \int_0^1x^r\F(x)^p(1-\F(x))^q\F'(x)\dd x
 &=Q_r(p+1,q+1).\label{eq:weightedFprime}
\end{align}
These identities are exact for complex parameters throughout their domains of
absolute convergence.

\section{A transform dictionary}

\subsection{Transforms of the inverse Fabius function}

The quantile substitution gives the master identity
\begin{equation}\label{eq:quantilepush}
 \int_0^1\Psi(\G(u))\dd u=\E\Psi(X).
\end{equation}
Therefore
\begin{align}
 \int_0^1e^{z\G(u)}\dd u
 &=M_X(z)
 =\prod_{j\ge1}\frac{e^{z/2^j}-1}{z/2^j},\label{eq:Gmgf}\\
 \int_0^1e^{-2\pi\ii\xi\G(u)}\dd u
 &=e^{-\pi\ii\xi}\Phi(\xi/2),\label{eq:Gfourier}\\
 \int_0^1(z-\G(u))^{-\nu}\dd u
 &=\E(z-X)^{-\nu},\label{eq:Gcauchy}\\
 \int_0^1\G(u)^r(1-\G(u))^s\dd u
 &=\E X^r(1-X)^s.\label{eq:Gmixed}
\end{align}
When \(s\in\N\), the last expression is the finite rational-moment sum
\begin{equation}\label{eq:Gmixedfinite}
 \sum_{k=0}^s(-1)^k\binom sk d_{r+k}.
\end{equation}

\subsection{Linear transforms of \(\F\)}

Integration by parts gives
\begin{align}
 \int_0^1e^{-sx}\F(x)\dd x
 &=\frac{M_X(-s)-e^{-s}}s,\label{eq:FLaplace}\\
 \widehat{\F\mathbf1_{[0,1]}}(\xi)
 &=\frac{e^{-\pi\ii\xi}\Phi(\xi/2)-e^{-2\pi\ii\xi}}
 {2\pi\ii\xi},\label{eq:FFourier}\\
 \int_0^1x^{s-1}\F(x)\dd x
 &=\frac{1-d_s}{s},\label{eq:FMellin}\\
 \int_0^1x^{s-1}(1-\F(x))\dd x
 &=\frac{d_s}{s}.\label{eq:survMellin}
\end{align}
The removable values at \(s=0\) or \(\xi=0\) are taken by continuity.

For \(\G\), one obtains the nonlinear Laplace dual
\begin{equation}\label{eq:GLaplace}
 \int_0^1e^{-su}\G(u)\dd u
 =\frac{\displaystyle\int_0^1e^{-s\F(x)}\dd x-e^{-s}}s
 =\frac{\displaystyle\sum_{n\ge0}\frac{(-s)^n}{n!}A_n-e^{-s}}s.
\end{equation}
Thus the Laplace transform of the quantile is encoded by the complete mixed-power
edge sequence \((A_n)\).

\subsection{Spacing transforms}

For \(p,q>0\), the spacing theorem yields an entire transform package.  With
\(N=p+q\),
\begin{align}
 \mathcal M_{p,q}(s)
 &:=\int_0^1x^{s-1}\F(x)^p(1-\F(x))^q\dd x\notag\\
 &=\frac{\E(Z_{p+1:N}^s-Z_{p:N}^s)}{s\binom Np},\label{eq:spacingMellin}\\
 \mathcal L_{p,q}(s)
 &:=\int_0^1e^{-sx}\F(x)^p(1-\F(x))^q\dd x\notag\\
 &=\frac{\E(e^{-sZ_{p:N}}-e^{-sZ_{p+1:N}})}{s\binom Np},\label{eq:spacingLaplace}\\
 \mathcal C_{p,q}(z)
 &:=\int_0^1\frac{\F(x)^p(1-\F(x))^q}{z-x}\dd x\notag\\
 &=\frac1{\binom Np}\E\log\frac{z-Z_{p:N}}{z-Z_{p+1:N}}.\label{eq:spacingCauchy}
\end{align}
More generally, for \(\nu\ne1\),
\begin{equation}\label{eq:spacingStieltjes}
 \int_0^1\frac{\F(x)^p(1-\F(x))^q}{(z-x)^\nu}\dd x
 =\frac{\E[(z-Z_{p+1:N})^{1-\nu}-(z-Z_{p:N})^{1-\nu}]}
 {(\nu-1)\binom Np}.
\end{equation}
All branch choices are understood consistently on a simply connected domain
avoiding \([0,1]\).

\section{The Cauchy--Stieltjes hierarchy of the \(\Up\)-law}

For \(z\in\C\setminus[-1,1]\) and \(\Re\nu>0\), define
\begin{equation}\label{eq:Snu}
 S_\nu(z)=\int_{-1}^1\frac{\Up(y)}{(z-y)^\nu}\dd y
 =\E(z-Y)^{-\nu}.
\end{equation}

\subsection{Laplace--sinh representation}

For \(\Re z>1\),
\begin{equation}\label{eq:SnuLaplace}
 S_\nu(z)=\frac1{\Gamma(\nu)}
 \int_0^\infty t^{\nu-1}e^{-zt}
 \prod_{j\ge1}\frac{\sinh(t/2^j)}{t/2^j}\dd t.
\end{equation}
Combining \eqref{eq:cumulants} with Watson expansion yields a complete
large-\(z\) Bell--Bernoulli asymptotic; because the support is compact, it is in
fact convergent for \(|z|>1\):
\begin{equation}\label{eq:Snumoments}
 S_\nu(z)=z^{-\nu}
 \sum_{m\ge0}\frac{(\nu)_{2m}}{(2m)!}\frac{c_m}{z^{2m}}.
\end{equation}
Every coefficient is rational.

\subsection{Dyadic functional equations}

\begin{theorem}[Dyadic Stieltjes hierarchy\tagnew]\label{thm:Snu}
For \(\nu\ne1\),
\begin{equation}\label{eq:Srec}
 \boxed{
 S_\nu(z)=\frac{2^{\nu-1}}{\nu-1}
 \left[S_{\nu-1}(2z-1)-S_{\nu-1}(2z+1)\right].}
\end{equation}
At \(\nu=1\),
\begin{equation}\label{eq:S1log}
 S_1(z)=\E\log\frac{2z+1-Y}{2z-1-Y}.
\end{equation}
Furthermore,
\begin{align}
 S_\nu^{(m)}(z)&=(-1)^m(\nu)_mS_{\nu+m}(z),\label{eq:Sderiv}\\
 S_1'(z)&=2\left[S_1(2z+1)-S_1(2z-1)\right],\label{eq:S1fde}\\
 S_1(-z)&=-S_1(z).\label{eq:Sodd}
\end{align}
\end{theorem}

\begin{proof}
The self-similarity \(Y=(V+Y')/2\), with
\(V\sim\operatorname{Unif}[-1,1]\) independent of \(Y'\), gives
\[
 S_\nu(z)=2^{\nu-1}\E\int_{-1}^1(2z-Y'-v)^{-\nu}\dd v.
\]
The elementary integral in \(v\) gives \eqref{eq:Srec}; its logarithmic limit is
\eqref{eq:S1log}.  Differentiation proves \eqref{eq:Sderiv}, and the case
\(\nu=2\) of \eqref{eq:Srec} gives \eqref{eq:S1fde}.  Symmetry of \(Y\) gives
\eqref{eq:Sodd}.
\end{proof}

\subsection{Boundary values and a Hilbert-transform refinement law}

For \(-1<x<1\), the upper boundary value satisfies
\begin{equation}\label{eq:Sboundary}
 S_{1,+}(x)=\PV\int_{-1}^1\frac{\Up(t)}{x-t}\dd t-\pi\ii\Up(x).
\end{equation}
Taking imaginary and real parts of \eqref{eq:S1fde} yields respectively the
known dyadic refinement of \(\Up\) and a parallel functional equation for its
Hilbert transform.  Thus the Cauchy transform packages the density and its
harmonic conjugate into a single dyadic object.

\subsection{Endpoint derivative growth}

At the right endpoint,
\begin{equation}\label{eq:Sendpointderiv}
 \frac{(-1)^n}{n!}S_1^{(n)}(1)
 =2^{-n-1}\E X^{-n-1}.
\end{equation}
The endpoint law \eqref{eq:endpointlog}, by a one-dimensional Laplace principle,
implies
\begin{equation}\label{eq:negmomentgrowth}
 \log\E X^{-s}=\frac{\log2}{2}s^2+O(s\log s).
\end{equation}
Hence
\begin{equation}\label{eq:Sdergrowth}
 \boxed{
 \lim_{n\to\infty}\frac1{n^2}
 \log\left|\frac{S_1^{(n)}(1)}{n!}\right|
 =\frac{\log2}{2}.}
\end{equation}
The one-sided derivatives at \(1\) all exist, but the Taylor series there has
radius zero.  This is a transform-level signature of the flat endpoint of
\(\Up\).

\section{Generalized resolvent duality for \(\F\) and \(\G\)}

For \(z\notin[0,1]\), define
\begin{align}
 C_{F,\nu}(z)&=\int_0^1\frac{\F(x)}{(z-x)^\nu}\dd x,\label{eq:CFnu}\\
 C_{G,\nu}(z)&=\int_0^1\frac{\G(u)}{(z-u)^\nu}\dd u.\label{eq:CGnu}
\end{align}

\begin{theorem}[CDF/quantile resolvent duality\tagnew]\label{thm:resdual}
For \(\nu\ne1\),
\begin{align}
 C_{F,\nu}(z)
 &=\frac{(z-1)^{1-\nu}-\E(z-X)^{1-\nu}}{\nu-1},\label{eq:CFdual}\\
 C_{G,\nu}(z)
 &=\frac{(z-1)^{1-\nu}
 -\displaystyle\int_0^1(z-\F(x))^{1-\nu}\dd x}{\nu-1}.
 \label{eq:CGdual}
\end{align}
At \(\nu=1\),
\begin{align}
 C_F(z)&=\E\log(z-X)-\log(z-1),\label{eq:CFlog}\\
 C_G(z)&=\int_0^1\log(z-\F(x))\dd x-\log(z-1).
 \label{eq:CGlog}
\end{align}
Both satisfy the same reflection law
\begin{equation}\label{eq:Creflection}
 C_F(z)-C_F(1-z)=C_G(z)-C_G(1-z)
 =\log\frac z{z-1}.
\end{equation}
\end{theorem}

\begin{proof}
Use the primitive
\((z-x)^{1-\nu}/(\nu-1)\) and integrate by parts.  In the quantile case,
\(\int H(u)\dd\G(u)=\int H(\F(x))\dd x\).  The logarithmic formulas are the
limits as \(\nu\to1\).  Reflection follows by substituting \(x\mapsto1-x\) and
using the symmetries of \(\F\) and \(\G\).
\end{proof}

Let
\begin{equation}\label{eq:Tdef}
 T(z)=\E(z-X)^{-1}=2S_1(2z-1).
\end{equation}
The affine self-similarity \(X=(U+X')/2\) yields
\begin{align}
 T(z)&=2\E\log\frac{2z-X}{2z-1-X},\label{eq:Tlog}\\
 T'(z)&=4[T(2z)-T(2z-1)].\label{eq:Tfde}
\end{align}
These formulas form a nonlinear/linear pair: \eqref{eq:Tlog} is a logarithmic
potential identity, while \eqref{eq:Tfde} is a linear dyadic
functional-differential equation.

\section{Sinc-product energies and the constant \(\int_0^1\F^2\)}

Set
\begin{equation}
 A_2=\int_0^1\F(x)^2\dd x.
\end{equation}
Numerically,
\begin{equation}\label{eq:A2numeric}
 A_2=0.404436767983985281465621491145\ldots.
\end{equation}
No familiar elementary or classical-special-function closed form is known in
the audited corpus.  There is, however, an exact sinc-product energy form.

\begin{theorem}[Sinc energy and complementary Gini sum rule\tagnew]
With \(\Phi\) as in \eqref{eq:phi},
\begin{align}
 A_2&=\int_0^\infty\Phi(\xi)^2\dd\xi,\label{eq:A2Phi}\\
 \E|X-X'|&=1-2A_2
 =\frac1{2\pi^2}\int_0^\infty
 \frac{1-\Phi(\xi)^2}{\xi^2}\dd\xi,\label{eq:GiniPhi}\\
 \boxed{
 \int_0^\infty\Phi(\xi)^2\dd\xi
 +\frac1{4\pi^2}\int_0^\infty
 \frac{1-\Phi(\xi)^2}{\xi^2}\dd\xi=\frac12.}
 \label{eq:sumrule}
\end{align}
\end{theorem}

\begin{proof}
Since \(\Up(x)=1-\F(x)\) on \([0,1]\), symmetry and Parseval give
\[
 \int_{\R}\Up(x)^2\dd x=2A_2
 =\int_{\R}\Phi(\xi)^2\dd\xi,
\]
which is \eqref{eq:A2Phi}.  The standard cycles-normalized fractional-energy
identity at exponent one is
\[
 \E|Z|=\frac1{\pi^2}\int_0^\infty
 \frac{1-\Re\varphi_Z(\xi)}{\xi^2}\dd\xi.
\]
For \(Z=X-X'\),
\(|\varphi_X(\xi)|^2=\Phi(\xi/2)^2\), yielding \eqref{eq:GiniPhi} after a
change of variables.  Finally,
\(\E|X-X'|=2\int_0^1\F(1-\F)=1-2A_2\).
\end{proof}

The integrand in \eqref{eq:GiniPhi} is regular at the origin:
\begin{equation}
 \lim_{\xi\to0}\frac{1-\Phi(\xi)^2}{\xi^2}=\frac{4\pi^2}{9},
\end{equation}
because \(\Var(Y)=1/9\).

For every integer \(q\ge2\), Fourier convolution gives the higher-power
representation
\begin{equation}\label{eq:AqFourier}
 A_q=\frac12\int_{\R^{q-1}}
 \Phi(\xi_1)\cdots\Phi(\xi_{q-1})
 \Phi\!\left(-\sum_{j=1}^{q-1}\xi_j\right)
 \dd\xi_1\cdots\dd\xi_{q-1}.
\end{equation}
For \(0<\alpha<2\), the full fractional energy is
\begin{equation}\label{eq:fractionalenergy}
 \E|X-X'|^\alpha
 =\frac{2^{1-\alpha}\Gamma(1+\alpha)\sin(\pi\alpha/2)}
 {\pi^{1+\alpha}}
 \int_0^\infty\frac{1-\Phi(\xi/2)^2}{\xi^{1+\alpha}}\dd\xi.
\end{equation}

\begin{conjecture}[Arithmetic opacity of the energy constant\tagconj]
The number \(A_2\) is transcendental.  More conservatively, it does not belong
to the field generated over \(\overline{\mathbb Q}\) by finitely many values of
Gamma, zeta, polylogarithm, and elementary functions at algebraic arguments.
\end{conjecture}

The conjecture is deliberately separated from the exact formulas
\eqref{eq:A2Phi} and \eqref{eq:sumrule}; those formulas may instead suggest a
new class of ``dyadic sinc periods.''

\section{General derivatives and antiderivatives of arbitrary order}

\subsection{Finite Thue--Morse derivative formulas}

Let \(s_2(j)\) be the binary digit sum and \(T_n=n(n+1)/2\).  Iterating
\eqref{eq:uprefinement} gives the baseline finite formulas
\begin{align}
 \Up^{(n)}(x)
 &=2^{T_n}\sum_{j=0}^{2^n-1}
 (-1)^{n+s_2(j)}
 \Up(2^nx+2j+1-2^n),\label{eq:upnth}\\
 \F^{(n)}(x)
 &=2^{T_n}\sum_{j=0}^{2^{n-1}-1}
 (-1)^{n-1+s_2(j)}
 \Up(2^nx+2j+1-2^n),\qquad n\ge1.
 \label{eq:Fnth}
\end{align}

For a mixed power
\(W_{p,q}(x)=\F(x)^p(1-\F(x))^q\), Fa\`a di Bruno gives the complete derivative
formula
\begin{equation}\label{eq:Wnth}
 W_{p,q}^{(n)}(x)
 =\sum_{k=0}^n
 \left.\frac{\dd^k}{\dd u^k}[u^p(1-u)^q]\right|_{u=\F(x)}
 B_{n,k}\bigl(\F'(x),\ldots,\F^{(n-k+1)}(x)\bigr),
\end{equation}
where \(B_{n,k}\) is a partial exponential Bell polynomial and
\begin{equation}\label{eq:outerderiv}
 \frac{\dd^k}{\dd u^k}[u^p(1-u)^q]
 =\sum_{r=0}^k\binom kr
 p^{\underline r}(-1)^{k-r}q^{\underline{k-r}}
 u^{p-r}(1-u)^{q-k+r}.
\end{equation}
Substituting \eqref{eq:Fnth} makes \eqref{eq:Wnth} a finite expression in
Thue--Morse translates of \(\Up\).

\subsection{Integer and fractional primitives of \(\Up\)}

The normalized \(n\)-fold primitive from \(-\infty\) is
\begin{equation}
 U_n(x)=\frac1{(n-1)!}\int_{-\infty}^x(x-t)^{n-1}\Up(t)\dd t.
\end{equation}
For \(x\le1\),
\begin{equation}\label{eq:upprimitive}
 U_n(x)=2^{\binom n2}\F\!\left(\frac{x+1}{2^n}\right),
\end{equation}
whereas for \(x\ge1\), compact support gives the moment polynomial
\begin{equation}\label{eq:upprimitivepoly}
 U_n(x)=\sum_{k=0}^{n-1}
 \frac{(-1)^k\mu_k}{k!(n-1-k)!}x^{n-1-k},
 \qquad \mu_k=\E Y^k.
\end{equation}
The odd \(\mu_k\) vanish, and the even ones are the Bell--Bernoulli moments in
\eqref{eq:cumulants}.

\subsection{Fractional primitives of \(\F\) and \(\G\)}

The baseline probabilistic form for \(\F\) is
\begin{equation}\label{eq:fracF}
 I^\nu\F(x)=\frac1{\Gamma(\nu+1)}\E(x-X)_+^\nu.
\end{equation}
For the inverse,
\begin{equation}\label{eq:fracG}
 I^\nu\G(y)=\frac1{\Gamma(\nu+1)}
 \int_0^{\G(y)}(y-\F(t))^\nu\dd t.
\end{equation}
The integer derivatives of the inverse are compactly expressed by the inverse
operator
\begin{equation}\label{eq:Gnth}
 \G^{(n)}(u)=
 \left[
 \left(\frac1{\F'(x)}\frac\dd{\dd x}\right)^{n-1}
 \frac1{\F'(x)}
 \right]_{x=\G(u)}.
\end{equation}
Equivalently, one may expand \eqref{eq:Gnth} using Bell polynomials in the
ratios \(\F^{(k)}/\F'\).  Equations \eqref{eq:Wnth}, \eqref{eq:fracspacing},
\eqref{eq:upprimitive}, and \eqref{eq:Gnth} together provide a uniform answer to
the ``\(n\)-th derivative/antiderivative'' part of the problem.

\section{Finite Thue--Morse splines and certified series approximations}

Let
\begin{equation}
 X_M=\sum_{j=1}^M\frac{U_j}{2^j},
 \qquad \F_M(x)=\Pp(X_M\le x).
\end{equation}
The generalized Irwin--Hall formula is
\begin{equation}\label{eq:FM}
 \F_M(x)=\frac{2^{M(M+1)/2}}{M!}
 \sum_{k=0}^{2^M-1}(-1)^{s_2(k)}
 \left(x-\frac{k}{2^M}\right)_+^M.
\end{equation}
Thus \(\F_M\) is a rational piecewise polynomial on the dyadic mesh.  The tail
coupling satisfies
\begin{equation}\label{eq:coupling}
 0\le X-X_M\le2^{-M}\qquad\text{almost surely}.
\end{equation}

\begin{theorem}[Certified rational approximation\tagnew]\label{thm:certified}
Let \(p,q\in\N_{>0}\), \(N=p+q\), and let \(H\) be Lipschitz on \([0,1]\) with
constant \(L\).  Define
\[
 J=\int_0^1H'(x)\F(x)^p(1-\F(x))^q\dd x,
 \qquad
 J_M=\int_0^1H'(x)\F_M(x)^p(1-\F_M(x))^q\dd x.
\]
Then
\begin{equation}\label{eq:certified}
 |J-J_M|\le\frac{2L}{\binom Np}\,2^{-M}.
\end{equation}
If \(H'\) is a polynomial with rational coefficients, then \(J_M\in\mathbb Q\)
and can be evaluated exactly by integrating finitely many dyadic polynomials.
\end{theorem}

\begin{proof}
Couple every observation with its truncation using the same uniforms.  By
\eqref{eq:coupling}, every corresponding order statistic differs by at most
\(2^{-M}\).  Apply \eqref{eq:spacing} to \(\F\) and \(\F_M\); the difference of
the two adjacent-gap increments of \(H\) is at most \(2L2^{-M}\).  Formula
\eqref{eq:FM} proves rationality.
\end{proof}

Consequently, every integer mixed moment has a rational sequence representation
with a certified exponential error.  This complements the more rapidly
convergent but structurally specialized \(q\)-Richardson schemes in the corpus.
For noninteger powers, one may combine \eqref{eq:FM} with the Abel series
\eqref{eq:abelseries}, or use the beta--quantile integral directly.

\section{Numerical experiments}

The archive contains \texttt{fabius\_integral\_transforms.py}.  It constructs a
uniform-grid fixed point for \(\F\), computes exact rational moments by
\eqref{eq:momentrec}, evaluates the sinc product, and checks independent sides
of the principal identities.  The complete output is in
\texttt{numerical\_results.txt}.  Selected results are shown below.

\begin{center}
\begin{tabular}{@{}ll@{}}
\toprule
Check & Absolute residual \\
\midrule
\(\F(1/4)=5/72\), grid \(2^{17}\) & \(1.29\times10^{-11}\) \\
Weighted spacing identity, \(e^x\F^2(1-\F)^3\) & \(3.35\times10^{-12}\) \\
Fractional spacing primitive, order \(1.65\) & \(1.26\times10^{-12}\) \\
Cauchy equation \eqref{eq:S1fde} & below \(2.3\times10^{-16}\) \\
Sinc/Gini sum rule \eqref{eq:sumrule} & \(8.1\times10^{-15}\) \\
\bottomrule
\end{tabular}
\end{center}

The two independent estimates
\[
 \int_0^1\F(x)^2\dd x\approx0.4044367679787557
\]
from the grid and
\[
 \int_0^{64}\Phi(\xi)^2\dd\xi
 \approx0.4044367679839853
\]
from oscillatory quadrature agree to the expected grid accuracy.  The normalized
balanced values begin
\[
 K_1\approx0.5733793921,
 \quad K_{10}\approx0.5030089983,
 \quad K_{100}\approx0.5000147741,
 \quad K_{1000}\approx0.50000000757,
\]
consistent with strict decrease and superalgebraic localization.  They are far
from the ultimate logarithmic regime in \eqref{eq:Klog}, as the Lambert saddle
moves only logarithmically.

\section{Conjectures and frontier directions}

\subsection{Analytic continuation and non-holonomicity}

\begin{conjecture}[No local crossing of the Cauchy cut\tagconj]
The Cauchy transform \(S_1\) admits no analytic continuation through any
nonempty open subinterval of \([-1,1]\).  Equivalently, the density \(\Up\) is
not real-analytic on any nonempty open interval.
\end{conjecture}

A proof would imply that \(S_1\) is not D-finite over \(\C(z)\), since a
D-finite function has only finitely many singular points in the finite plane.
The dyadic equation \eqref{eq:S1fde}, together with the nowhere-analytic
structure of \(\F\), is the natural route.

\subsection{The mixed-power resolvent as a new special function}

The function \(R\) in \eqref{eq:Rdef} is a Stieltjes-type transform of the
symmetric random variable \(\F(U)\), where \(U\sim\operatorname{Unif}[0,1]\).
It obeys the involutive Möbius law \eqref{eq:Rmobius}, but no closed form is
known.

\begin{problem}
Determine the maximal analytic continuation of \(R\), its boundary values on
\([1,\infty)\), and whether its Möbius functional equation plus positivity
characterizes the Fabius pushforward measure.
\end{problem}

\begin{problem}
Develop Pad\'e and continued-fraction approximants to \(R\) from the exact
Hausdorff moments \(A_n\), and compare their pole distributions with the
\(q\)-Richardson accelerants already present in the corpus.
\end{problem}

\subsection{Balanced Lambert tomography}

The full conjecture \cref{conj:Kfull} offers a new way to probe the periodic
endpoint correction: balanced beta localization samples the flat central germ
of \(\G'\), which is itself a reparameterized copy of the endpoint germ of
\(\F\).  Numerically fitting \(K_m\) after removing the Lambert saddle may
recover the endpoint phase without evaluating \(\F\) at extremely small
arguments.

\begin{problem}
Derive the first two explicit amplitude and phase terms in
\eqref{eq:Kfullconj}, and design a confluent \(q\)-Richardson scheme in the
variable \(w_m=W(16m\log2)\) that preserves the periodic phase rather than
averaging it away.
\end{problem}

\subsection{Dyadic sinc periods}

The constants
\begin{equation}
 \mathcal P_q=\int_{\R^{q-1}}
 \Phi(\xi_1)\cdots\Phi(\xi_{q-1})
 \Phi\!\left(-\sum\xi_j\right)\dd\bm\xi=2A_q
\end{equation}
form a natural period family attached to the dyadic sinc product.

\begin{problem}
Find recurrences, functional equations, or arithmetic obstructions for
\(\mathcal P_q\).  Determine whether the complementary energy identity
\eqref{eq:sumrule} is the first member of a hierarchy involving fractional
Riesz energies and higher convolution powers.
\end{problem}

\subsection{Formalization roadmap}

The most Lean-friendly new results are:
\begin{enumerate}
 \item the random-interval identity and adjacent-spacing theorem, using finite
       exchangeability and order statistics;
 \item the Riemann--Liouville spacing formula, by Tonelli and an elementary
       truncated-power integral;
 \item beta--quantile Pascal/reflection identities;
 \item the Cauchy hierarchy, from the affine fixed-point law and integration of
       a uniform variable;
 \item the sinc-energy identity, from Parseval and the standard Fourier formula
       for \(|x|\);
 \item the certified finite-spline bound, from the deterministic tail coupling.
\end{enumerate}
The balanced logarithmic asymptotic requires the endpoint theorem plus a
quantitative Laplace principle and should follow only after the elementary
transform layer is formalized.

\section{Compact formula compendium}

For convenient reference, the central formulas are collected here.

\begin{longtable}{>{\raggedright\arraybackslash}p{0.29\textwidth}p{0.66\textwidth}}
\toprule
Object & Formula \\
\midrule
\endfirsthead
\toprule
Object & Formula \\
\midrule
\endhead
Weighted mixed integral &
\(\displaystyle
\int h\F^p(1-\F)^q
=\binom{p+q}{p}^{-1}\E[H(Z_{p+1})-H(Z_p)]\).\\[0.75em]
Fractional mixed primitive &
\(\displaystyle
I^\nu[\F^p(1-\F)^q](x)
=\frac{\E[(x-Z_p)_+^\nu-(x-Z_{p+1})_+^\nu]}
{\Gamma(\nu+1)\binom{p+q}{p}}\).\\[0.75em]
Incomplete monomial primitive &
Equation \eqref{eq:incompleteQ} in terms of truncated \(Q_{m+1}\).\\[0.5em]
Beta--quantile reflection &
\(Q_1(a,b)+Q_1(b,a)=B(a,b)\).\\[0.5em]
Mixed-power recurrence &
\(I_{p,q}=I_{p+1,q}+I_{p,q+1}=(-1)^q\Delta^qA_p\).\\[0.5em]
Bivariate generating function &
\(\displaystyle\sum I_{p,q}z^pw^q=[zR(z)+wR(w)]/(z+w-zw)\).\\[0.75em]
Balanced normalization &
\(K_m=(2m+1)\binom{2m}{m}I_{m,m}\downarrow1/2\), with
\(\log(K_m-1/2)\sim-(\log m)^2/(8\log2)\).\\[0.75em]
Power of \(\Up\) &
\(\displaystyle\int_{-1}^1|x|^\alpha\Up(x)^q\dd x
=\frac{2q}{\alpha+1}Q_{\alpha+1}(1,q)\).\\[0.75em]
Inverse Fourier transform &
\(\displaystyle\int_0^1e^{-2\pi\ii\xi\G(u)}\dd u
=e^{-\pi\ii\xi}\Phi(\xi/2)\).\\[0.75em]
\(\Up\)-Stieltjes hierarchy &
\(\displaystyle S_\nu(z)=\frac{2^{\nu-1}}{\nu-1}
[S_{\nu-1}(2z-1)-S_{\nu-1}(2z+1)]\).\\[0.75em]
CDF resolvent &
\(\displaystyle C_{F,\nu}(z)=
[(z-1)^{1-\nu}-\E(z-X)^{1-\nu}]/(\nu-1)\).\\[0.75em]
Sinc energy &
\(\displaystyle A_2=\int_0^\infty\Phi^2\), and equation
\eqref{eq:sumrule}.\\[0.5em]
Finite spline &
\(\displaystyle\F_M(x)=2^{M(M+1)/2}M!^{-1}
\sum_{k<2^M}(-1)^{s_2(k)}(x-k/2^M)_+^M\).\\
\bottomrule
\end{longtable}

\section{Conclusion}

The integration theory of the Fabius--Rvachev system becomes substantially more
coherent when it is organized around three dual viewpoints:
\begin{enumerate}
 \item \emph{random intervals}: mixed CDF powers are occupation probabilities
       of gaps between two independent sample blocks;
 \item \emph{beta quantiles}: the substitution \(u=\F(x)\) turns mixed powers
       into beta localization against \(\G'\);
 \item \emph{self-similar transforms}: the affine fixed-point laws of \(X\) and
       \(Y\) induce exact dyadic equations for Cauchy, Stieltjes, Laplace, and
       Fourier transforms.
\end{enumerate}
These viewpoints cover indefinite integrals, all integer and fractional
primitives, broad families of definite transforms, and nonlinear inverse-Fabius
integrals without requiring a fictitious elementary closed form.  They also
expose new frontier phenomena: strict sign-regularity, balanced
beyond-all-orders localization, endpoint derivative growth, and dyadic sinc
periods.  The resulting formulas are compatible with the repository's
q-binomial, Lambert, Bell, Bernoulli, Thue--Morse, and sinc-product structures,
but are not duplications of those earlier developments.

\begin{thebibliography}{99}

\bibitem{AriasDefinition}
J. Arias de Reyna,
\emph{An infinitely differentiable function with compact support: Definition
and properties}, arXiv:1702.05442, 2017.

\bibitem{AriasArithmetic}
J. Arias de Reyna,
\emph{Arithmetic of the Fabius function}, arXiv:1702.06487, 2017.

\bibitem{Fabius1966}
J. Fabius,
\emph{A probabilistic example of a nowhere analytic \(C^\infty\)-function},
Zeitschrift f\"ur Wahrscheinlichkeitstheorie und Verwandte Gebiete
\textbf{5} (1966), 173--174.

\bibitem{Haugland}
J. K. Haugland,
\emph{Evaluating the Fabius function}, arXiv:1609.07999, 2016/2020.

\bibitem{Have}
S. G. Have,
\emph{A recurrence relation for the odd order moments of the Fabius function},
arXiv:1703.01904, 2017.

\bibitem{Rvachev}
V. L. Rvachev and V. A. Rvachev,
\emph{Non-classical Methods of Approximation Theory in Boundary Value Problems},
Naukova Dumka, Kiev, 1979 (in Russian).

\bibitem{ProveIt}
V. Reshetnikov,
\emph{ProveIt: Analysis/FabiusFunction documentation and formal development},
GitHub repository, audited at commit
\path{a6555a64671b54f851563a04a391bb7845bd6571}, 2026.

\bibitem{DavidNagaraja}
H. A. David and H. N. Nagaraja,
\emph{Order Statistics}, 3rd ed., Wiley, 2003.

\bibitem{Widder}
D. V. Widder,
\emph{The Laplace Transform}, Princeton University Press, 1941.

\end{thebibliography}

\end{document}
